\section{ATOMIC 与 COMET：Language Model 不等于 Knowledge Model}

这一章从评价转向表示。通用语言模型优化连续文本，而 commonsense system 往往需要回答 typed question：某个事件的原因是什么、参与者想要什么、接下来会发生什么。ATOMIC 用 symbolic relation 定义问题空间，COMET 用 neural generation 扩展答案。

\subsection{Symbolic Graph 与 Neural Model 的分工}

slide 把 ATOMIC 和 COMET 并排：前者由人类 crowdsourcing 写成 symbolic commonsense knowledge graph，后者在这些 relation-conditioned triples 上训练成为 neural commonsense model。二者不是竞争关系；graph 提供 schema 与监督，model 提供泛化和生成。

\lecturefigure{slide-13-language-models-not-knowledge-models.jpg}{Language model 与 knowledge model 的目标不相同}{官方录像恢复幻灯片 V013；Stanford CS25 Lecture 16}

下一页加入 COMET-ATOMIC 2020，强调 machine-readable graph 与 generative model 的循环：人类先定义 relation ontology 与 seed knowledge，模型学习 $p(t\mid h,r)$，再为新事件生成 tail inference。这里 $h$ 是 head event，$r$ 是关系类型，$t$ 是生成的 commonsense conclusion。

\lecturefigure{slide-14-atomic-comet-human-symbolic-neural.jpg}{ATOMIC、COMET 与 COMET-ATOMIC 2020 的 symbolic-neural 关系}{官方录像恢复幻灯片 V014；Hwang et al., 2021}

形式上，COMET 的训练目标可写为
\begin{equation}
\mathcal{L}_{\mathrm{COMET}}(\theta)
=-\mathbb{E}_{(h,r,t)\sim\mathcal{D}}
\log p_{\theta}(t\mid h,r).
\end{equation}
$\mathcal{D}$ 是知识三元组数据，$h$ 与 $r$ 明确限定要生成哪类推断。与 unrestricted continuation 相比，这个 interface 更接近“查询知识模型”。

\subsection{Relation Types 是知识的坐标系}

完整 ATOMIC 图把 social interaction、physical entity、event-centered 等 relation types 展开。读图时不要逐节点背诵，而要看颜色和边的语义：同一事件可以连接 agent intent、agent reaction、other reaction、precondition 与 effect。relation type 让“可能的推断”变成可分解任务，也让 evaluator 能分别检查不同能力。

\lecturefigure{slide-15-atomic2020-full-graph.jpg}{ATOMIC 2020：事件节点与 23 类常识关系的知识图谱}{官方录像恢复幻灯片 V015；Hwang et al., 2021}

\begin{knowledgebox}{Typed Relation 为什么重要}
若只让模型自由续写，“合理”可以有无数方向。指定 \texttt{xIntent}、\texttt{xEffect}、\texttt{oReact} 或 physical relation 后，数据收集、训练与错误分析都获得坐标。代价是 ontology 本身决定系统能表达什么，未建模的关系不会自动出现。
\end{knowledgebox}

\subsection{COMET、GPT-3 与 GPT-2 的比较边界}

比较 slide 给出 human-judged commonsense quality：COMET-BART、few-shot GPT-3 与 zero-shot GPT-2 的设置不同。最先应读的是方法标签而不是数字高低；COMET 使用专门数据与监督，GPT-3 依赖 few-shot prompting，GPT-2 更小且 zero-shot。结果支持“专门知识模型可以更小、更好用”，但不能推出 architecture 单因素的因果结论。

\lecturefigure{slide-16-comet-vs-gpt3-gpt2.jpg}{COMET 与通用语言模型的 human-judged commonsense 对比}{官方录像恢复幻灯片 V016；Stanford CS25 Lecture 16}

\teachervoice{讲者主动说这不是 apple-to-apple comparison。她同时承认从 GPT-2 到 GPT-3 的 scale jump 很有趣，又指出 GPT-3 对许多工程师和研究者太大。课堂的重点是 system utility：一个更小、可查询、能嵌入下游系统的模型可能更有价值。}

\subsection{Knowledge Backbone 的复用价值}

最后一页把 COMET / ATOMIC 放在中心，周围连接 persona-aware conversation、figurative language、storytelling、fantasy games 与 interactive learning。它展示的不是所有应用都被“解决”，而是统一 relation interface 降低了重复构建常识组件的成本。应用性能仍受 coverage、quality 与 cultural scope 限制。

\lecturefigure{slide-17-commonsense-backbone-applications.jpg}{ATOMIC / COMET 作为多个下游任务的 commonsense backbone}{官方录像恢复幻灯片 V017；Stanford CS25 Lecture 16}

\subsection{本章小结}

ATOMIC 把常识问题写成 typed symbolic queries，COMET 把图谱监督压进 generative model。语言模型与知识模型的差别不在是否使用 Transformer，而在训练目标、接口和验收方式。下一章进一步问：如果人工写知识图谱太慢，能否让大模型生产语料，再用 critic 把它蒸馏成更小模型？

